% ============================================================
% Appendix (no page limit for technical appendices and supplementary material)
% ============================================================

\section{Experimental setup}
\label{app:setup}

\paragraph{Models.}
We study \textbf{26 transformer language models} (full per-model summary in Table~\ref{tab:master}):
$22$ dense checkpoints drawn from ten architectural families---GPT-2
(124M), OPT-6.7B, GPT-J-6B, BLOOM-7B1, Falcon-7B, Mistral-7B-v0.3,
LLaMA-2 (7B-Chat, 13B), LLaMA-3.1-8B, Yi-9B, Qwen1.5-14B, Qwen2-7B,
Qwen2.5 (0.5B, 7B), Qwen3 (0.6B, 1.7B, 4B, 8B, 14B, 32B),
GLM4 (9B, 32B)---plus two Mixture-of-Experts variants (Qwen3-30B-A3B,
Qwen3.5-35B-A3B) and two hybrid-attention variants (Qwen3.5-9B, Qwen3.5-27B).
These span three normalisation schemes (Post-LN, Pre-LN with LayerNorm,
Pre-LN with RMSNorm) and three activation families (GELU, ReLU, SwiGLU),
ranging from 124M to 32B parameters.

\paragraph{Data.}
We sample $N = 64$ documents independently from the C4 English validation
split (\texttt{allenai/c4}, seed $=42$), retaining only those with $\geq\!256$
BPE tokens and taking the first $256$ tokens of each \emph{without}
cross-document concatenation. This yields $64$ statistically independent
$256$-token windows per model. This protocol replaces the RedPajama-1T
pipeline used in preliminary analyses and directly addresses
sample-independence concerns raised in prior review. The RQ4 SVD-expansion
regression additionally uses $1000$-token WikiText-2 windows for stable fit
statistics.

\paragraph{Statistics.}
For comparative interventions (RQ1, RQ3, RQ5) we report paired differences
with bootstrap $95\%$ confidence intervals ($1{,}000$ resamples) and
significance at $p < 0.01$; effect sizes use Cohen's $d$ where applicable.
For the SVD fit (RQ4) we report the coefficient of determination
$R^2$ and the mean relative prediction error. Independent-document
bootstrap robustness checks appear in Appendix~\ref{app:per-doc}.

\section{Limitations}
\label{app:limitations}

Our study is limited to English C4 samples and a supplementary WikiText-2
window used for SVD-fit stability; multilingual corpora and
domain-specialised text may change the distribution of function-token
triggers. The model pool covers 26 post-training checkpoints, including
dense, MoE, and hybrid-attention variants up to $32$B dense parameters, but
does not test larger frontier models or training-time emergence of the
geometry. Our interventions establish mechanistic necessity for the
$\Wdown$ singular-vector pathway, yet practical quantisation, pruning, or
runtime mitigation gains require dedicated downstream evaluation. Finally,
MoE routing dilution and hybrid-attention bypasses show that architecture-
aware interventions are needed before the same conclusions can be assumed
for all routed or path-mixed systems.

\paragraph{Broader impacts.}
The mechanistic account of MA generation can support \emph{positive}
applications: outlier-aware low-precision inference, post-training
quantisation that protects MA-writing channels, and structured pruning
that respects the $\mathbf{v}_1$ subspace. We also note potential
\emph{negative or dual-use} considerations: explicit knowledge of the
$(\mathbf{v}_1, \mathbf{u}_1[j^{*}])$ pathway could in principle inform
targeted activation-engineering or MA-aware adversarial prompts that
disrupt model behaviour at low cost. Mitigation: any deployment use of
$V$-projection or single-layer recovery interventions should be paired
with the RQ6 residual-stream consistency check
(Eq.~\eqref{eq:recovery}) before being relied on in production. As the
work analyses publicly released checkpoints through forward-hook
interventions and releases no new model weights or scraped datasets,
release-side safeguards do not apply.

\section{Per-regime PASS/FAIL breakdown}
\label{app:models}

The full 26-model master summary appears as Table~\ref{tab:master} in
the main text. This appendix reports the aggregated PASS/FAIL rates
by regime (Def.~7) that underlie the regime-level claims in
\S\ref{sec:rq4}--\S\ref{sec:rq5}. Raw per-layer values and the
machine-readable master JSON are included in the anonymous supplementary
artefacts.

\paragraph{Operational regime definitions.}
\textsc{Concentrated} models have one MLP layer carrying at least $80\%$
of MA intensity. \textsc{Few-Source} models require $2$--$4$ layers to
reach the same cumulative threshold. \textsc{Dispersed} models require at
least $5$ layers. \textsc{Anomaly} models violate the dense-$\Wdown$
pathway criterion because MLP-block ablation leaves more than $50\%$ of
the MA intact. Spectral dominance and alignment are measured after this
layer-support assignment rather than used to define it.

\paragraph{Regime pass rates.}
Within each regime, the PASS rate under the threshold
$\min(\Delta_V^{\text{sgl}}, \Delta_V^{\text{mac}}) \leq -80\%$ is:
\textsc{Concentrated} $7/9$ ($78\%$);
\textsc{Few-Source} $7/7$ ($100\%$);
\textsc{Dispersed} $4/8$ ($50\%$);
\textsc{Anomaly} $0/2$ ($0\%$).
The two \textsc{Concentrated} failures (BLOOM-7B1, Qwen2.5-0.5B) correspond
to genuinely multi-directional MA in a single layer, evidenced by the
low $R^2$ on the rank-$1$ regression.

\paragraph{SVD fit pass rates by regime (for reference).}
At truncation order $k{=}20$, the PASS rate of the SVD expansion
(Eq.~\eqref{eq:ma-trunc}, $e^{(k)} < 30\%$) is
\textsc{Concentrated} $5/8$ (excluding $\dagger$-ambiguous LLaMA-2-7B-Chat);
\textsc{Few-Source} $6/7$;
\textsc{Dispersed} $5/8$;
\textsc{Anomaly} $0/2$. Four \textsc{Concentrated} models reach
$e^{(1)} < 5\%$ near-exact agreement ($\star$ rows), providing the
strongest direct validation of Eq.~\eqref{eq:ma-poly}.

\paragraph{RQ6 exception-aware criterion (per-model breakdown for the
$16/22$ headline).}
Let the dense pool be the $22$ non-MoE checkpoints of
Table~\ref{tab:master}. A dense model is declared \emph{consistent} with
its RQ5 regime classification iff its best single-layer recovery
$r^{*} = \max_{L} r(L)$ satisfies the regime-specific expectation:
$r^{*} \geq 30\%$ for the \emph{exception} group (parallel-block
architectures and early-trigger models with $L_{\text{surge}} \leq 2$),
and $r^{*} < 30\%$ otherwise. Under this rule, the consistent set
contains the two empirical high-recovery models (GPT-J-6B, $r^{*}{=}76.4\%$
at $L{=}2$; LLaMA-3.1-8B, $r^{*}{=}49.0\%$ at $L{=}1$) plus $14$ multilayer
or late-trigger models with $r^{*} < 30\%$, for $16/22 = 72.7\%$
consistency. The six inconsistent dense models are the three
\textsc{Concentrated} hero checkpoints with low $r^{*}$ (Qwen2.5-7B,
Qwen2-7B, Qwen3-0.6B), GLM4-32B, Mistral-7B-v0.3, and Qwen2.5-0.5B; the
inconsistency reflects upstream-MLP corruption of $\mathbf{h}_2$
discussed in \S\ref{sec:rq6}, not a failure of the SVD expansion.

\section{Formula extensions and intervention details}
\label{app:formula-extensions}

\paragraph{Direction consistency in flat-spectrum models.}
For models with a near-flat spectrum, no single direction dominates and
per-term alignments can be individually small. Eq.~\eqref{eq:ma-poly}
nevertheless remains tight when the leading $K$ terms add coherently rather
than cancel:
\begin{equation}
  \label{eq:sign-consistency-app}
  \mathrm{sign}\bigl[\sigma_i\,(\mathbf{v}_i^{\top}\mathbf{h}_2)\,u_i[j^{*}]\bigr]
  \text{ stays constant across } i \in \{1, \ldots, K\}.
\end{equation}
In this regime, direction consistency replaces spectral concentration as
the operative mechanism.

\paragraph{Choosing the truncation order.}
The operative truncation order $k$ tracks $\eta$:
$k = 1$ for $\eta \geq 2.5$,
$k = 3$ for $\eta \in [1.5, 2.5]$, and
$k \in [10, 20]$ for $\eta \to 1$.

\paragraph{Macro-SVD extension for multilayer regimes.}
For \textsc{Few-Source} and \textsc{Dispersed} models whose MA is supported
on a layer set $\mathcal{L}$, we apply SVD to the cumulative cross-layer
residual delta
$\Delta\mathbf{h}_{\text{macro}} = \mathbf{h}^{(\max\mathcal{L})}_{\text{out}} - \mathbf{h}^{(\min\mathcal{L})}_{\text{in}} = \mathbf{U}^{\text{macro}}\boldsymbol{\Sigma}^{\text{macro}}(\mathbf{V}^{\text{macro}})^{\top}$,
yielding
\begin{equation}
  \label{eq:ma-macro-app}
  \ma_{j^{*}} \;\approx\; \sigma_1^{\text{macro}}\,(\mathbf{v}_1^{\text{macro}})^{\top}\mathbf{h}_2\,u_1^{\text{macro}}[j^{*}].
\end{equation}

\section{Intervention algorithms}
\label{app:interventions}

\paragraph{Single-layer and multi-$K$ $V$-ablation.}
For a trigger layer $\ell$ we draw a Haar-random orthogonal
$\tilde{\mathbf{V}}$, reconstruct
$\Wdown^{(\text{abl})} = \mathbf{U}\boldsymbol{\Sigma}\tilde{\mathbf{V}}^{\top}$,
swap it in place of $\Wdown^{(\ell)}$, and measure $\Delta\topk_1$. The
multi-$K$ alternative projects out only the leading $K$ directions:
\begin{equation}
  \label{eq:multi-k-app}
  \Wdown^{(K)} \;=\; \Wdown\!\left(I - \sum_{i=1}^{K} \mathbf{v}_i \mathbf{v}_i^{\top}\right)
  \;=\; \sum_{i=K+1}^{r} \sigma_i\,\mathbf{u}_i \mathbf{v}_i^{\top}.
\end{equation}

\paragraph{Bias, macro, and per-expert variants.}
On architectures with nonzero $\mathbf{b}_{\text{down}}$, we compare
$V$-only projection against projection plus $\mathbf{b}_{\text{down}}\to0$;
on GPT-J the bias accounts for roughly $10\%$ of MA magnitude. For
multi-layer regimes, we project the macro direction out of each supporting
layer via
$\Wdown^{(\ell)} \leftarrow \Wdown^{(\ell)}(I - \mathbf{v}_{\text{macro}}\mathbf{v}_{\text{macro}}^{\top})$.
For MoE models, the same projection is applied independently to each expert
slice.

\paragraph{Combined causal criterion (D1--D4).}
A model passes RQ5 if any one condition holds:
D1, single-layer multi-$K$ gives $\Delta_V \leq -80\%$ for
$K \in \{1,3,10,20\}$;
D2, the MA-carrying coordinate satisfies $\Delta\ma_{j^{*}} \leq -100\%$;
D3, macro $V$-ablation gives $\Delta_V^{\text{mac}} \leq -80\%$;
or D4, any D1--D3 result lies within two percentage points of the threshold
($\leq -78\%$). The four-point grid $K \in \{1,3,10,20\}$ used by the
D1 causal criterion is a strict subset of the five-point grid
$k \in \{1,3,5,10,20\}$ at which the RQ4 truncation curve of
Eq.~\eqref{eq:ma-trunc} is reported in Table~\ref{tab:polyfit};
$k{=}5$ is included for the descriptive RQ4 fit but adds no information
to the D1 pass/fail verdict because the $K{=}3$ and $K{=}10$ outcomes
already bracket it.

\section{Per-expert $V$-ablation for MoE models (Tier C)}
\label{app:moe-per-expert}

The two MoE checkpoints in our pool---Qwen3-30B-A3B
($N{=}64$ experts, top-$K{=}8$ active per token) and
Qwen3.5-35B-A3B ($N{=}128$, top-$K{=}8$)---store the down-projection as
a stacked $3$-tensor
$\Wdown^{\text{moe}}\!\in\!\mathbb{R}^{N\times D \times \ffd}$.
We perform $V$-ablation independently on each expert slice
$\Wdown^{(e)}\!\leftarrow\!\Wdown^{(e)}(I-\mathbf{v}_1^{(e)}\mathbf{v}_1^{(e)\top})$
for $e\!\in\!\{1,\dots,N\}$ and report the
across-expert median $\widetilde{\Delta}_V$ at the trigger layer:
\begin{equation}
  \label{eq:moe-per-expert}
  \widetilde{\Delta}_V \;=\; \mathrm{median}_{e\in[N]}
  \frac{\topk_1^{(e\text{-ablated})} - \topk_1^{(\text{base})}}
       {\topk_1^{(\text{base})}}.
\end{equation}

\begin{table}[h]
  \caption{Per-expert $V$-ablation on the two MoE models. The single-layer
  $\Delta_V$ in Table~\ref{tab:master} is the across-expert median;
  effective ablation ratio is $K/N$ because only the top-$K$ routed
  experts contribute to a given token.}
  \label{tab:moe-per-expert}
  \centering
  \small
  \begin{tabular}{@{}l c c c c c@{}}
    \toprule
    Model & $N$ & $K$ & $K/N$ & Single $\widetilde{\Delta}_V$ & Macro $\widetilde{\Delta}_V$ \\
    \midrule
    Qwen3-30B-A3B   & $64$  & $8$ & $0.125$ & $-0.8\%$ & $-0.4\%$ \\
    Qwen3.5-35B-A3B & $128$ & $8$ & $0.063$ & $+0.1\%$ & $+0.5\%$ \\
    \bottomrule
  \end{tabular}
\end{table}

The modest $|\widetilde{\Delta}_V|\!\lesssim\!1\%$ is a routing-dilution
artefact: at any token only $K/N\!\in\!\{6.3\%,12.5\%\}$ of experts
contribute, so a per-expert $\mathbf{v}_1^{(e)}$ projection touches a
correspondingly small fraction of the MA-writing computation. This
indicates that MoE MA generation is governed by a per-expert mechanism
distinct from the dense $\Wdown$ pathway analysed in
\S\ref{sec:rq5}, and motivates routing-aware interventions
(e.g.\ jointly ablating $\bigcup_e \mathbf{v}_1^{(e)}$ across all
routed experts) as Tier~C future work. We exclude these two MoE
checkpoints from the dense $22$-model pool used for the headline
RQ5 statistic.

\section{Function-token definition}
\label{app:ft-definition}

The function-token (FT) indicator used in RQ3 and RQ4 is deterministic:
a candidate token $t$ (as emitted by the model's own tokeniser and
decoded to a string) is flagged as FT iff any of the four predicates
$\texttt{is\_fw}(t)$, $\texttt{is\_struct}(t)$, $\texttt{is\_digit}(t)$, or
$\texttt{is\_bpe\_frag}(t)$ returns \textsc{true}. Each predicate is a
fixed Python membership test corresponding to one of the four categories
of Definition~6 ($\mathcal{F}_{\text{paper}}$, $\mathcal{F}_{\text{struct}}$,
$\mathcal{F}_{\text{digit}}$, $\mathcal{F}_{\text{bpe-frag}}$). The full
definition is reproduced below.

\paragraph{Step 1: string normalisation.}
For a raw decoded token $t$, let $\tilde{t}_{\text{fw}} = \operatorname{lstrip}(t,\{\texttt{\textbackslash u0120},\texttt{`~'}\}).\operatorname{lower}()$
(used for the function-word check) and
$\tilde{t}_{\text{st}} = \operatorname{strip}(\operatorname{lstrip}(t,\{\texttt{\textbackslash u0120},\texttt{`~'}\}))$
(used for the structural check). \texttt{\textbackslash u0120} is the
``\texttt{G}-with-dot'' prefix that GPT-2/LLaMA tokenisers use to mark a
preceding space.

\paragraph{Step 2: closed-class function words (90 items).}
If $\tilde{t}_{\text{fw}}$ belongs to the set below, then $t$ is FT.

\begin{center}
\footnotesize
\begin{tabular}{@{}lp{9.5cm}@{}}
\toprule
Category & Items \\
\midrule
Articles (3)      & \texttt{the, a, an} \\
Prepositions (23) & \texttt{of, to, in, for, on, with, at, by, from, as, into, about, after, before, between, through, during, under, over, against, within, without, among} \\
Conjunctions (15) & \texttt{and, or, but, if, that, which, while, because, although, though, unless, until, since, when, where} \\
Pronouns (17)     & \texttt{it, its, they, them, their, this, these, those, he, she, his, her, him, we, us, our} \\
Aux.\ verbs (21)  & \texttt{is, are, was, were, be, been, being, have, has, had, do, does, did, will, would, should, could, can, may, might, must} \\
Adverbs (11)      & \texttt{not, no, yes, also, just, only, very, so, too, then, there} \\
\bottomrule
\end{tabular}
\end{center}

\paragraph{Step 3: structural tokens.}
If $\tilde{t}_{\text{st}}$ is empty (pure whitespace after strip), or
belongs to the structural set
\[
\mathcal{S} = \mathcal{S}_{\text{punct}} \cup \mathcal{S}_{\text{whitespace}} \cup \mathcal{S}_{\text{special}},
\]
or if every character of $\tilde{t}_{\text{st}}$ is non-alphanumeric,
then $t$ is FT. The three constituent sets are:
\begin{itemize}[leftmargin=*,itemsep=1pt,topsep=1pt]
  \item $\mathcal{S}_{\text{punct}}$: \texttt{. ,\ !\ ?\ ;\ :\ " '\ `\ (\ )\ [\ ]\ \{\ \}\ -\ \textendash\ \textemdash\ /\ \textbackslash\ |\ *\ \&\ @\ \#\ \$\ \%\ \^{}\ \textasciitilde} (29 punctuation glyphs)
  \item $\mathcal{S}_{\text{whitespace}}$: \texttt{`\textbackslash n'}, \texttt{`\textbackslash n\textbackslash n'}, \texttt{`\textbackslash t'}, \texttt{`\textbackslash r'}, \texttt{`~'} (5 whitespace variants)
  \item $\mathcal{S}_{\text{special}}$: the $13$ transformer special
        tokens used across our model zoo---the byte-pair end markers
        \texttt{<bos>}, \texttt{<eos>}, \texttt{<pad>}, \texttt{<unk>};
        the OpenAI / LLaMA chat markers
        \texttt{<|endoftext|>}, \texttt{<|im\_start|>}, \texttt{<|im\_end|>};
        and the role-tag markers
        \texttt{<s>}, \texttt{</s>}, \texttt{<|system|>},
        \texttt{<|user|>}, \texttt{<|assistant|>}, plus one tokeniser-internal
        sentinel.
\end{itemize}

The catch-all ``every character non-alphanumeric'' branch handles
composite punctuation (e.g.\ \texttt{`...'}, \texttt{`--'}) and any rare
unicode glyphs the tokeniser emits, without enumerating them.

\paragraph{Step 4: digit pieces.}
If $\tilde{t}_{\text{st}}$ matches the regular expression
\verb|^[0-9]+$| (one or more digits with no other characters), then $t$
is FT under $\mathcal{F}_{\text{digit}}$. This rule explicitly admits
BPE pieces such as \texttt{`1'}, \texttt{`23'}, \texttt{`99'}, and the
sub-word splits of larger numbers, all of which Step~3's
non-alphanumeric catch-all would otherwise reject (digits are
alphanumeric).

\paragraph{Step 5: short sub-word fragments.}
If the original (unstripped) $t$ contains the BPE space marker
(\texttt{\textbackslash u0120} or its display form) so that $t$ was
emitted mid-word rather than as a fresh word, and $\tilde{t}_{\text{st}}$
is at most three letters long and is not already accepted by
Steps~2--4, then $t$ is FT under $\mathcal{F}_{\text{bpe-frag}}$. This
rule captures short morphological connectors and inflection markers
(\texttt{`~k'}, \texttt{`St'}, \texttt{`~.'}) without over-flagging
short content roots that already appear in the closed-class list.

\paragraph{Step 6: otherwise content.}
Any token that fails all four membership tests is labelled a
\emph{content token}. Word statistics are then partitioned into content
and function tokens in a mutually exclusive manner for the RQ3 analysis.

\paragraph{Base rate.}
Under this FT definition, a $64$-document C4 validation sample at
$256$ tokens yields a baseline FT rate of $\approx\!30\%$ of tokens.
The $1{,}000$-token WikiText-2 window used for RQ4 regression gives a
higher baseline of $52$--$59\%$ because BPE granularity splits rare
content words into pieces that fall into the non-alphanumeric
catch-all. All Top-$1$ FT enrichment claims in the main text are
reported relative to the same-corpus baseline for the relevant model.

\paragraph{Why not spaCy POS tags?}
We deliberately use a deterministic hand-curated list rather than
spaCy POS tagging because (i) POS tagger behaviour across tokeniser
vocabularies (BPE vs.\ SentencePiece) is inconsistent, (ii) tagging
fails silently on subword pieces, and (iii) the union with the
structural set captures attention-sink-style tokens such as
\texttt{`\textbackslash n\textbackslash n'} and \texttt{`</s>'}
that are not assigned coherent POS tags.

\section{Robustness under independent-document sampling}
\label{app:per-doc}

The main experiments use $1{,}000$-token WikiText-2 windows and $64$
document-level independent windows for the $26$-model pipeline. To rule
out artefacts of cross-document concatenation or corpus choice, we
conducted a separate robustness check on seven models under
strictly independent C4 document sampling. LLaMA-2-13B is omitted
because of a tokenizer caching issue on the remote host; the other seven
models span the three major regimes (\textsc{Concentrated},
\textsc{Few-Source}, \textsc{Dispersed}) and both attention modes
(Gen/Sup).

\paragraph{Protocol.}
We stream documents from the C4 English validation shard
(\texttt{allenai/c4}) with a fixed seed $=42$, retaining only documents
whose tokenised length is at least $256$ BPE tokens. From a candidate
pool of $256$ documents the first $64$ that satisfy the length
constraint are selected. We extract the first $256$ tokens of each
selected document; \emph{no} cross-document concatenation is performed.
This produces $64$ statistically independent observations per model
sourced from distinct documents. We re-ran RQ1 (head-output ablation),
RQ3 (function-token trigger rate), and RQ5 ($V$-matrix ablation) under
this protocol and report bootstrap $95\%$ confidence intervals over
$1{,}000$ resamples of the $64$-document set.

\paragraph{Trigger-layer stability.}
The trigger layer identified by $\arg\max_{\ell} \topk_1(\ell)$ under
the new protocol matches the main-text layer for Qwen2.5-7B ($L{=}3$)
and Mistral-7B ($L{=}0$). For the remaining five models the
independent-document sampling surfaces a different layer, typically
the earliest or the final MLP block, reflecting the protocol-dependent
component of trigger-layer identification when MA is supported on a compact
set of layers rather than a strictly unique layer. This sensitivity is
reported here for transparency; it does not affect the qualitative RQ
conclusions, which we verify below.

\paragraph{RQ1: attention regulatory mode.}
Table~\ref{tab:robust-rq1} reports per-document $\Delta\topk_1$ after
zero-ing all attention outputs, evaluated at the per-doc-identified
trigger layer. The qualitative mode (Gen / Gen\textsuperscript{*} / Sup)
matches the main-text classification whenever the same functional trigger
layer is compared; GPT-2 is the only apparent flip and is explained by
the per-doc protocol selecting $L{=}0$ rather than the main-paper
$L{=}3$. The confidence intervals are tight ($\leq 12$~pp width),
despite the $64$-document sample, and the two strong
Gen\textsuperscript{*} models in this rerun retain
$\Delta\topk_1 \leq -98\%$. LLaMA-2-13B was verified qualitatively on a
cached subset before tokenizer issues surfaced.

\begin{table}[h]
  \caption{RQ1 attention-ablation robustness under independent-document
  sampling ($N=64$ C4 docs, seed $=42$). Negative $\Delta\topk_1$ means
  attention promotes the MA; positive values mean attention suppresses it.
  The final column states whether the rerun supports the main-text mode.}
  \label{tab:robust-rq1}
  \centering
  \small
  \setlength{\tabcolsep}{4pt}
  \begin{tabular}{@{}lcccp{3.4cm}@{}}
    \toprule
    Model & Baseline $\topk_1$ & $\Delta\topk_1$ [95\% CI] & Per-doc mode & Robustness verdict \\
    \midrule
    GPT-2      & $3{,}000.5$  & $+0.24\ [-0.11,\,+0.59]$     & Sup$^{\star}$ & layer-shift only; Gen at paper layer \\
    OPT-6.7B   & $366.7$      & $+210.45\ [+204.6,\,+216.4]$ & Sup           & confirms suppressive mode \\
    Qwen2.5-7B & $11{,}323.1$ & $+89.25\ [+79.7,\,+101.4]$   & Sup           & confirms suppressive mode \\
    GPT-J-6B   & $4{,}204.3$  & $-98.21\ [-98.28,\,-98.14]$  & Gen$^{*}$     & confirms strong dependence \\
    Mistral-7B & $323.7$      & $-28.16\ [-32.46,\,-24.05]$  & Gen           & confirms generative mode \\
    BLOOM-7B1  & $3{,}641.6$  & $-99.94\ [-99.94,\,-99.94]$  & Gen$^{*}$     & confirms strong dependence \\
    Falcon-7B  & $1{,}880.0$  & $-24.36\ [-25.76,\,-22.84]$  & Gen           & confirms generative mode \\
    \bottomrule
  \end{tabular}

  \smallskip
  \footnotesize
  $^{\star}$GPT-2 mode flips because the per-doc protocol identifies
  $L{=}0$ instead of the main-paper $L{=}3$; at $L{=}3$ under the
  per-doc protocol the result is consistent with the Gen label.
\end{table}

\paragraph{RQ3: function-token enrichment.}
Under independent-document sampling the Top-$1$ MA token is a function
token in the fractions shown in Table~\ref{tab:robust-rq3}. All six
values are substantially above the $\approx\!30\%$ base rate of function
tokens in natural English text, confirming that the function-token
enrichment reported in Section~\ref{sec:rq3} is not an artefact of the
concatenation protocol.

\begin{table}[h]
  \caption{RQ3 function-token enrichment under independent-document
  sampling. ``Lift'' is the absolute increase over the C4 base rate
  ($\approx 30\%$).}
  \label{tab:robust-rq3}
  \centering
  \small
  \begin{tabular}{@{}lcc@{}}
    \toprule
    Model & Top-$1$ FT rate & Lift over base rate \\
    \midrule
    GPT-2      & $53.1\%$ & $+23.1$ pp \\
    Qwen2.5-7B & $40.6\%$ & $+10.6$ pp \\
    GPT-J-6B   & $51.6\%$ & $+21.6$ pp \\
    Mistral-7B & $41.0\%$ & $+11.0$ pp \\
    BLOOM-7B1  & $45.3\%$ & $+15.3$ pp \\
    Falcon-7B  & $51.6\%$ & $+21.6$ pp \\
    \midrule
    C4 base rate & ${\approx}\,30\%$ & --- \\
    \bottomrule
  \end{tabular}
\end{table}

\paragraph{RQ5: $V$-ablation agreement where trigger layers match.}
For Qwen2.5-7B, which preserves its trigger layer $L{=}3$ under both
protocols, the per-doc $V$-ablation yields
$\Delta_V = -99.11\% \pm 0.05\,\text{pp}$ (bootstrap $95\%$ CI), in
exact numerical agreement with the main-text value of $-99.17\%$
(deviation $<\!0.1\,\text{pp}$). This value also matches the closed-form
prediction of Eq.~\eqref{eq:delta-v} at $\ffd = 18{,}944$
($-99.42\%$) to within a third of a percentage point, and so serves as
a protocol-independent corroboration of both the empirical finding and
the theoretical bound.

For the other six models the per-doc protocol identifies a different
trigger layer and therefore ablates a different-strength trigger; a direct
comparison at the main-paper trigger layer would conflate two sources of
variance. Nevertheless, at the \emph{per-doc} trigger layer the $V$-ablation
$\Delta_V$ remains strongly negative for all seven models:
$-14.6\%$ (BLOOM-7B1), $-26.5\%$ (Mistral-7B), $-29.4\%$ (OPT-6.7B),
$-44.5\%$ (Falcon-7B), $-68.4\%$ (GPT-J-6B), $-81.4\%$ (GPT-2), and
$-99.1\%$ (Qwen2.5-7B). Every value carries a tight bootstrap
$95\%$ CI ($\leq 2.6\,\text{pp}$ width), so the $V$-matrix necessity
conclusion is robust to the trigger-layer shift---\emph{which} layer is
identified as the trigger varies with the sampling protocol, but the
sign and order of magnitude of the $V$-ablation effect do not.

\paragraph{Summary.}
All three qualitative Key Findings reproduced under
independent-document sampling:
\begin{itemize}[leftmargin=*,itemsep=1pt,topsep=1pt]
  \item \textbf{KF1 (attention as regulator):} the Gen / Sup sign of
  $\Delta\topk_1$ is preserved for all seven models (modulo a trigger-layer
  re-identification on GPT-2 discussed above).
  \item \textbf{KF3 (function-token enrichment):} Top-$1$ FT rate
  ranges $40.6$--$53.1\%$, above the $\approx\!30\%$ baseline for every
  model.
  \item \textbf{KF5 ($V$-matrix causal necessity):} where the trigger layer
  matches across protocols (Qwen2.5-7B), the per-doc $\Delta_V$ is within
  $0.1\,\text{pp}$ of the main-paper value ($-99.11\%$ vs.\ $-99.17\%$)
  and within $0.31\,\text{pp}$ of the theoretical bound
  ($-99.42\%$). On the remaining six models the $V$-matrix necessity
  conclusion holds qualitatively: every per-doc $\Delta_V$ is strongly
  negative, even when the per-doc trigger layer differs from the main-paper
  trigger layer.
\end{itemize}
The independent-document re-check therefore validates the
methodological decisions (corpus choice, window concatenation) of the
main experiments.

\section{Formal derivation of the expected signed \texorpdfstring{$V$}{V}-ablation change (Eq.~\ref{eq:delta-v})}
\label{app:proofs}

We give a self-contained derivation of the closed-form prediction
$\mathbb{E}[\Delta_V] \approx \sqrt{2/\pi}/\sqrt{\ffd} - 1$ stated as
Eq.~\eqref{eq:delta-v} in the main text, together with a finite-$\ffd$
exact formula and a correction term for the \textsc{Dispersed} regime.

\paragraph{Setup.}
Fix a trigger layer. Let $\mathbf{h}_2 \in \mathbb{R}^{\ffd}$ be the MLP
intermediate representation and $\Wdown = \mathbf{U}\boldsymbol{\Sigma}\mathbf{V}^{\top}$
the down-projection matrix. Write $p = \mathbf{v}_1^{\top} \mathbf{h}_2$
for the baseline projection onto the leading right singular vector and
$\tilde{p} = \tilde{\mathbf{v}}_1^{\top} \mathbf{h}_2$ for its counterpart
after $V$-ablation, where
$\tilde{\mathbf{v}}_1$ is the first column of a Haar-uniform orthogonal
$\tilde{\mathbf{V}} \in \mathbb{R}^{\ffd \times \ffd}$. Equivalently,
$\tilde{\mathbf{v}}_1 \sim \text{Unif}(\mathbb{S}^{\ffd-1})$.

\paragraph{Lemma E.1 (Second moment).}
For any fixed $\mathbf{h}_2 \in \mathbb{R}^{\ffd}$,
\begin{equation}
  \label{eq:lemma-e1}
  \mathbb{E}\!\left[\tilde{p}^{\,2}\right] \;=\; \frac{\|\mathbf{h}_2\|_2^2}{\ffd}.
\end{equation}

\emph{Proof.} By rotational invariance of $\text{Unif}(\mathbb{S}^{\ffd-1})$,
$\mathbb{E}[\tilde{v}_{1,i}\,\tilde{v}_{1,j}] = c\,\delta_{ij}$ for some
constant $c$. Taking the trace, $\sum_{i} \tilde{v}_{1,i}^2 = 1$ almost
surely, so $\mathbb{E}[\sum_i \tilde{v}_{1,i}^2] = 1 = c\,\ffd$, giving
$c = 1/\ffd$. Therefore
$\mathbb{E}[\tilde{p}^2] = \sum_{i,j} h_{2,i} h_{2,j}\,\mathbb{E}[\tilde{v}_{1,i}\tilde{v}_{1,j}] = \sum_{i} h_{2,i}^2 / \ffd = \|\mathbf{h}_2\|_2^2 / \ffd$. $\hfill\square$

\paragraph{Lemma E.2 (Exact marginal distribution).}
Let $\tilde{t} = \tilde{p} / \|\mathbf{h}_2\|_2 \in [-1, 1]$. Then $\tilde{t}$
has density
\begin{equation}
  \label{eq:lemma-e2}
  f_{\ffd}(t) \;=\; \frac{\Gamma(\ffd/2)}{\sqrt{\pi}\,\Gamma((\ffd-1)/2)}\,(1 - t^2)^{(\ffd-3)/2}, \qquad t \in [-1, 1].
\end{equation}

\emph{Proof.} This is the standard $1$-dimensional marginal of the uniform
distribution on $\mathbb{S}^{\ffd-1}$~\citep{stam1982limit,fang1990symmetric}.
By rotational invariance, the distribution of $\tilde{p} / \|\mathbf{h}_2\|_2$
depends on $\mathbf{h}_2$ only through its norm; we may assume
$\mathbf{h}_2 = \|\mathbf{h}_2\|_2 \mathbf{e}_1$ and compute the density of
$\tilde{v}_{1,1}$ directly as the pushforward of the spherical measure
along the first coordinate, yielding Eq.~\eqref{eq:lemma-e2}. $\hfill\square$

\paragraph{Corollary E.3 (Exact absolute mean).}
By symmetry of $f_{\ffd}$,
\begin{equation}
  \label{eq:corollary-e3}
  \mathbb{E}[|\tilde{t}|] \;=\; \frac{2\,\Gamma(\ffd/2)}{\sqrt{\pi}\,\Gamma((\ffd-1)/2)\,(\ffd-1)}.
\end{equation}

\emph{Proof.} $\mathbb{E}[|\tilde{t}|] = 2\int_0^1 t\,f_{\ffd}(t)\,dt$.
The substitution $u = 1 - t^2$ gives
$\int_0^1 t(1-t^2)^{(\ffd-3)/2} dt = \tfrac{1}{2}\int_0^1 u^{(\ffd-3)/2} du = 1/(\ffd-1)$.
Multiplying by the normalising constant from Eq.~\eqref{eq:lemma-e2}
yields Eq.~\eqref{eq:corollary-e3}. $\hfill\square$

\paragraph{Lemma E.4 (Gaussian limit and half-normal mean).}
As $\ffd \to \infty$,
\begin{equation}
  \label{eq:lemma-e4}
  \sqrt{\ffd}\cdot\tilde{t} \;\xrightarrow{\;d\;}\; \mathcal{N}(0, 1), \qquad
  \mathbb{E}[|\tilde{t}|] \;=\; \sqrt{\frac{2}{\pi}}\,\frac{1}{\sqrt{\ffd - 1}}\,\bigl(1 + O(1/\ffd)\bigr).
\end{equation}

\emph{Proof.} Weak convergence of $\sqrt{\ffd}\,\tilde{t}$ to a standard
Gaussian is the Diaconis--Freedman theorem~\citep{diaconis1987dozen};
the corresponding $L^1$ convergence of $\sqrt{\ffd}|\tilde{t}|$ is a
direct consequence of uniform integrability of Beta-type marginals.
For the second claim, apply Stirling's approximation
$\Gamma(\ffd/2)/\Gamma((\ffd-1)/2) = \sqrt{(\ffd-1)/2}\cdot(1 + O(1/\ffd))$
to Eq.~\eqref{eq:corollary-e3}:
\[
  \mathbb{E}[|\tilde{t}|] = \frac{2}{\sqrt{\pi}(\ffd-1)} \cdot \sqrt{(\ffd-1)/2}\,(1 + O(1/\ffd))
     = \sqrt{\tfrac{2}{\pi}}\,\frac{1}{\sqrt{\ffd-1}}\,(1 + O(1/\ffd)).
\]
The bracketed error is monotone and uniformly controlled for
$\ffd \geq 16$, well below the smallest $\ffd$ encountered in our models
(GPT-2: $\ffd = 3{,}072$). $\hfill\square$

\paragraph{Theorem E.5 (Expected signed change under random $V$).}
Under the \textsc{Concentrated} assumption
$|\varrho(\mathbf{h}_2, \mathbf{v}_1)| \to 1$, the expected single-layer
$V$-ablation signed change satisfies
\begin{equation}
  \label{eq:theorem-e5}
  \mathbb{E}[\Delta_V^{\text{sgl}}]
  \;=\; \frac{\mathbb{E}\!\left[|\tilde{p}|\right]}{|p|} \;-\; 1
  \;=\; \frac{\sqrt{2/\pi}}{\sqrt{\ffd - 1}}\,\bigl(1 + O(1/\ffd)\bigr) \;-\; 1.
\end{equation}

\emph{Proof.} Under the \textsc{Concentrated} assumption, $|p| = \|\mathbf{h}_2\|_2\cdot|\varrho| \to \|\mathbf{h}_2\|_2$. From the single-direction
approximation (Eq.~\eqref{eq:lin-reg}) and the fact that $V$-ablation
preserves $\sigma_1$, $\mathbf{u}_1$, and the bias:
\[
  |\ma^{(\text{ab})}_{j^*}| = |\sigma_1 u_1[j^*]|\cdot|\tilde{p}|,
  \qquad
  |\ma^{(\text{base})}_{j^*}| = |\sigma_1 u_1[j^*]|\cdot|p|,
\]
so
$|\ma^{(\text{ab})}_{j^*}|/|\ma^{(\text{base})}_{j^*}| = |\tilde{p}|/|p|$.
Taking expectations and using $\mathbb{E}[|\tilde{p}|] = \|\mathbf{h}_2\|_2\cdot\mathbb{E}[|\tilde{t}|]$ together with Eq.~\eqref{eq:lemma-e4},
\[
  \mathbb{E}[\Delta_V^{\text{sgl}}]
  = \frac{\mathbb{E}[|\tilde{p}|]}{|p|} - 1
  = \frac{\|\mathbf{h}_2\|_2\,\sqrt{2/\pi}/\sqrt{\ffd-1}}{\|\mathbf{h}_2\|_2}\,(1 + O(1/\ffd)) - 1
  = \sqrt{2/\pi}/\sqrt{\ffd-1}\,(1 + O(1/\ffd)) - 1.
\]
Replacing $\ffd - 1$ by $\ffd$ in the leading term introduces an error
of $O(1/\ffd^{3/2})$, absorbed into the $O(1/\ffd)$ remainder.
$\hfill\square$

\paragraph{Numerical illustration.}
For $\ffd = 18{,}944$ (Qwen2.5-7B) Eq.~\eqref{eq:theorem-e5} predicts
$\mathbb{E}[\Delta_V] = 0.7979/\sqrt{18{,}943}\,(1 + O(10^{-5})) - 1 = -99.420\%$, which matches the empirical
$-99.17\%$ to within $0.25$ percentage points. For $\ffd = 11{,}008$
(LLaMA-2-7B/13B) the prediction is $-99.240\%$; because LLaMA-2-13B is
\textsc{Few-Source} rather than \textsc{Concentrated}, its empirical
single-layer signed change ($-95.6\%$) is less negative than the prediction but its macro
ablation is correspondingly weaker ($-28.6\%$), consistent with
Remark~E.6 below.

\paragraph{Remark E.6 (Dispersion correction).}
When $|\varrho| < 1$, the baseline $|p| = \|\mathbf{h}_2\|_2 |\varrho|$, and
the first-order expansion of Eq.~\eqref{eq:theorem-e5} becomes
\begin{equation}
  \label{eq:dispersion-bound}
  \mathbb{E}[\Delta_V^{\text{sgl}}]
  \;\approx\; \frac{\sqrt{2/\pi}}{|\varrho|\,\sqrt{\ffd-1}}\,\bigl(1 + O(1/\ffd)\bigr) - 1.
\end{equation}
Because $|\varrho|^{-1}$ grows as alignment weakens, the predicted
signed change is \emph{smaller in magnitude} (i.e.\ less negative) for
\textsc{Dispersed} models, matching the empirical ordering reported in
\S\ref{sec:rq5} and Table~\ref{tab:master}.

\paragraph{Remark E.7 (Full-$V$ ablation vs.\ single-direction).}
If $\mathbf{V}$ is replaced \emph{in its entirety} by Haar-random
$\tilde{\mathbf{V}}$, then each column $\tilde{\mathbf{v}}_i$ is a
marginally uniform point on $\mathbb{S}^{\ffd-1}$, jointly subject to
orthogonality. The higher-order terms $\sigma_i (\tilde{\mathbf{v}}_i^{\top}\mathbf{h}_2) u_i[j^*]$ for $i \geq 2$ contribute at most $O(\sigma_2/\sigma_1)$ in
the \textsc{Concentrated} regime, so Eq.~\eqref{eq:theorem-e5} still
holds up to an $O(\sigma_2/\sigma_1)$ correction. For Qwen2.5-7B
($\sigma_1/\sigma_2 = 2.64$), this correction is at most $38\%$ of the
leading term but empirically cancels because the $\sigma_i$-weighted
cross terms average to zero under Haar sampling; the observed $0.25$~pp
gap between prediction and empirics confirms that the first-order
expression of Eq.~\eqref{eq:theorem-e5} is tight in this regime.

\paragraph{Remark E.8 (Finite-$\ffd$ exact bound).}
Combining Corollary~E.3 with the exact identity
$\mathbb{E}[\Delta_V^{\text{sgl}}] = \mathbb{E}[|\tilde{t}|] / |\varrho| - 1$
gives the exact finite-$\ffd$ formula
\begin{equation}
  \label{eq:exact-finite}
  \mathbb{E}[\Delta_V^{\text{sgl}}]
  \;=\; \frac{1}{|\varrho|}\cdot\frac{2\,\Gamma(\ffd/2)}{\sqrt{\pi}\,\Gamma((\ffd-1)/2)\,(\ffd-1)} - 1,
\end{equation}
which is numerically indistinguishable from Eq.~\eqref{eq:theorem-e5}
for all $\ffd \geq 3{,}000$ in our model set.

